%! Confidence Intervals for the Difference of Two Means — Unit 7, Lesson 7.6 — Lesson Plan
\documentclass[10pt]{article}
\usepackage{apstats-article}
\usepackage{apstats-boxes}
\usepackage{graphicx}
\graphicspath{{images/}}
\newcommand{\TallMath}[1]{$\displaystyle #1\rule[-1.4em]{0pt}{3.2em}$}
\newcommand{\UnitNumberName}{Unit 7: Inference for Quantitative Data: Means \quad}
\newcommand{\LessonNumberName}{Lesson 7.6: Confidence Intervals for the Difference of Two Means}

\begin{document}
\begin{center}
    {\Large\bfseries \CourseName: \SchoolYear \\
    {\small\bfseries \UnitNumberName \LessonNumberName}}
\end{center}

\begin{tcolorbox}[colback=sky, colframe=navy, boxrule=0.9pt, arc=2mm,
    left=2mm, right=2mm, top=1.5mm, bottom=1.5mm]
    \textbf{Primary Objective:} Students will identify the two-sample $t$-interval as the
    appropriate procedure for estimating a difference of two population means, verify the
    required conditions, compute the standard error and margin of error, and calculate and
    interpret the resulting confidence interval.
    (Big Idea: UNC)
\end{tcolorbox}

\renewcommand{\arraystretch}{1.6}
\begin{skillbox}[Priority Ideas \& Skills]{goldbox}
\begin{minipage}[t]{0.48\linewidth}
    \textbf{AP Skill 1 --- Selecting Statistical Methods}
    \begin{itemize}
        \item Identify the two-sample $t$-interval as the appropriate procedure for a
              difference of two population means (Skill 1.D; UNC-4.V).
    \end{itemize}
    \textbf{AP Skill 4 --- Statistical Argumentation}
    \begin{itemize}
        \item Verify conditions for the two-sample $t$-interval: independence (random
              samples or randomized experiment; 10\% rule if sampling without replacement)
              and Normal (both $n_1$ and $n_2 \ge 30$ if skewed) (Skill 4.C; UNC-4.W).
    \end{itemize}
\end{minipage}
\hfill
\begin{minipage}[t]{0.48\linewidth}
    \textbf{AP Skill 3 --- Using Probability and Simulation}
    \begin{itemize}
        \item Determine the margin of error: $ME = t^* \cdot SE$ where
              $SE = \sqrt{s_1^2/n_1 + s_2^2/n_2}$ (Skill 3.D; UNC-4.X).
        \item Calculate the confidence interval:
              $(\bar x_1 - \bar x_2) \pm t^*\sqrt{s_1^2/n_1 + s_2^2/n_2}$
              with $df$ from technology (Skill 3.D; UNC-4.Y).
    \end{itemize}
    \textbf{Key Understandings}
    \begin{itemize}
        \item Degrees of freedom come from technology (not $n_1 + n_2 - 2$).
        \item An interval that does not contain 0 suggests the population means differ.
    \end{itemize}
\end{minipage}
\end{skillbox}

\begin{skillbox}[{Vocabulary, Concepts, \& Theorems}]{greenbox}
\begin{minipage}[t]{0.48\linewidth}
    \begin{itemize}
        \item \textbf{Two-sample $t$-interval for $\mu_1 - \mu_2$} --- a confidence
              interval procedure for the difference of two independent population means
              when $\sigma_1$ and $\sigma_2$ are unknown
        \item \textbf{Point estimate of $\mu_1 - \mu_2$} --- $\bar x_1 - \bar x_2$,
              the difference of sample means
        \item \textbf{Standard error of $\bar x_1 - \bar x_2$} ---
              $SE = \sqrt{\dfrac{s_1^2}{n_1} + \dfrac{s_2^2}{n_2}}$;
              estimates the standard deviation of the sampling distribution
    \end{itemize}
\end{minipage}
\hfill
\begin{minipage}[t]{0.48\linewidth}
    \begin{itemize}
        \item \textbf{Degrees of freedom (two-sample)} --- found using technology;
              not simply $n_1 + n_2 - 2$; always use calculator or software
        \item \textbf{Margin of error (ME)} --- $t^* \cdot SE$; determines the width
              of the interval around the point estimate
        \item \textbf{Critical value $t^*$} --- the $t$-value for the desired confidence
              level, based on $df$ from technology
        \item \textbf{Independence condition} --- two independent random samples; 10\%
              rule applies to each sample separately
    \end{itemize}
\end{minipage}
\end{skillbox}

\begin{skillbox}[Activate Prior Knowledge \& Spiral Review]{sky}
\begin{minipage}[t]{0.55\linewidth}\vspace{0pt}
    \textbf{Spiral Review (warm-up)}
    \begin{itemize}
        \item One-sample $t$-interval: $\bar x \pm t^* \cdot (s/\sqrt{n})$ --- Lesson 7.2
        \item SE formula and how it changes with sample size --- Lesson 7.2/7.3
        \item Conditions for one-sample inference: random and Normal --- Lesson 7.2
        \item Interpreting a confidence interval: ``We are $C\%$ confident the true
              mean is between \ldots'' --- Lesson 7.3
    \end{itemize}
\end{minipage}
\hfill
\begin{minipage}[t]{0.38\linewidth}\vspace{0pt}\centering
    % Authored warm-up compiles to target/ — no source PDF to embed here.
\end{minipage}
\end{skillbox}

\begin{skillbox}[Hook]{sky}
\begin{minipage}[t]{0.55\linewidth}
    \textbf{City A vs.\ City B: Which City Pays More? (3--4 min)}\\
    Display two data summaries: City A ($n_1 = 32$ workers, $\bar x_1 = \$62{,}400$,
    $s_1 = \$8{,}500$) and City B ($n_2 = 28$ workers, $\bar x_2 = \$58{,}700$,
    $s_2 = \$7{,}200$). Ask: ``Is the \$3,700 difference meaningful, or just random
    sampling variation?''

    Key discussion points:
    \begin{itemize}
        \item We have \emph{two} independent samples and want to estimate $\mu_1 - \mu_2$.
        \item Neither $\sigma$ is known, so we use a $t$-interval (not $z$).
        \item We need to account for \emph{both} samples' variability in the SE.
        \item An interval that doesn't include 0 suggests the means genuinely differ.
    \end{itemize}
\end{minipage}
\hfill
\begin{minipage}[t]{0.40\linewidth}
    \textbf{Bridge from Lessons 7.2--7.3}\\
    One-sample $t$-interval: $\bar x \pm t^* \cdot \dfrac{s}{\sqrt{n}}$\\[6pt]
    Two-sample extension:
    $(\bar x_1 - \bar x_2) \pm t^* \cdot \sqrt{\dfrac{s_1^2}{n_1} + \dfrac{s_2^2}{n_2}}$\\[6pt]
    Same logic; add the two sources of variability under the radical.
\end{minipage}
\end{skillbox}

\begin{skillbox}[Lesson: Sampling Distribution \& Choosing the Method]{sky}
\begin{minipage}[t]{0.48\linewidth}
    \textbf{Sampling Distribution of $\bar x_1 - \bar x_2$}
    \begin{itemize}
        \item Two independent SRSs from populations with means $\mu_1$, $\mu_2$ and
              standard deviations $\sigma_1$, $\sigma_2$.
        \item If populations are Normal \emph{or} both $n_1, n_2 > 30$:
              $\bar x_1 - \bar x_2$ is approximately Normal with
              mean $\mu_1 - \mu_2$ and
              $SD = \sqrt{\dfrac{\sigma_1^2}{n_1} + \dfrac{\sigma_2^2}{n_2}}$
        \item Since $\sigma_1, \sigma_2$ are unknown, replace with $s_1, s_2$ and
              use a $t$-distribution.
    \end{itemize}
\end{minipage}
\hfill
\begin{minipage}[t]{0.48\linewidth}
    \textbf{Appropriate Procedure (UNC-4.V.2)}\\
    Use a \textbf{two-sample $t$-interval for a difference of population means} when:
    \begin{enumerate}
        \item Two independent samples (not matched pairs).
        \item Both population standard deviations are unknown.
        \item Conditions are met (see next box).
    \end{enumerate}
    \textit{Never pool variances on the AP Exam unless the problem explicitly states
    equal population variances.}
\end{minipage}
\end{skillbox}

\begin{skillbox}[Lesson (cont.): Conditions, SE, and the Interval]{sky}
\begin{minipage}[t]{0.48\linewidth}
    \textbf{Conditions to Verify (UNC-4.W)}
    \begin{enumerate}
        \item \textbf{Independence:}
              \begin{itemize}
                  \item Each sample is a random sample or part of a randomized experiment.
                  \item If sampling without replacement: $n_1 \le 10\%N_1$ and
                        $n_2 \le 10\%N_2$.
                  \item The two groups must be independent of each other.
              \end{itemize}
        \item \textbf{Normal:}
              \begin{itemize}
                  \item If both $n_1, n_2 \ge 30$: CLT applies; condition is met.
                  \item If either $n < 30$: check that sample is free of strong skew
                        and outliers.
              \end{itemize}
    \end{enumerate}
\end{minipage}
\hfill
\begin{minipage}[t]{0.48\linewidth}
    \textbf{The Confidence Interval (UNC-4.X/Y)}
    \[
        (\bar x_1 - \bar x_2) \pm t^* \sqrt{\frac{s_1^2}{n_1} + \frac{s_2^2}{n_2}}
    \]
    \begin{itemize}
        \item $\bar x_1 - \bar x_2$: point estimate of $\mu_1 - \mu_2$.
        \item $SE = \sqrt{s_1^2/n_1 + s_2^2/n_2}$: pooled standard error.
        \item $t^*$: critical value; $df$ from technology (TI-84: 2-SampTInt).
        \item $ME = t^* \cdot SE$.
    \end{itemize}
    \textit{Interpretation:} ``We are $C\%$ confident the true difference $\mu_1 - \mu_2$
    is between \textit{lower} and \textit{upper}.''
\end{minipage}
\end{skillbox}

\begin{skillbox}[Active Monitoring]{redbox}
\begin{minipage}[t]{0.48\linewidth}
    \textbf{Circulate and Check}
    \begin{itemize}
        \item Students use $SE = \sqrt{s_1^2/n_1 + s_2^2/n_2}$, not a pooled formula.
        \item Conditions verified for \emph{both} samples separately.
        \item $df$ from technology, not $n_1 + n_2 - 2$.
        \item Interval interpretation references $\mu_1 - \mu_2$ (not $\bar x_1 - \bar x_2$)
              and includes context.
        \item Students identify whether 0 is in the interval and state the implication.
    \end{itemize}
\end{minipage}
\hfill
\begin{minipage}[t]{0.48\linewidth}
    \textbf{Cold-Call Prompts}
    \begin{itemize}
        \item ``Why do we add $s_1^2/n_1$ and $s_2^2/n_2$ under the radical?''
        \item ``If the interval is $(−2.1, 8.3)$, what does that tell us about $\mu_1 - \mu_2$?''
        \item ``If the interval is entirely positive, what does that tell us?''
        \item ``What would change if we switched Group 1 and Group 2?''
    \end{itemize}
\end{minipage}
\end{skillbox}

\begin{skillbox}[Group Work \& Differentiation]{redbox}
\begin{multicols}{3}
    \textbf{Tier R --- Remediate}
    \begin{itemize}
        \item Provide the SE formula and CI template with labeled blanks.
        \item Focus on identifying conditions and substituting values.
        \item Activity: Tier R scenario only.
    \end{itemize}
    \columnbreak
    \textbf{Tier A --- Approaching}
    \begin{itemize}
        \item Complete Tier R and Tier A scenarios independently.
        \item Justify each condition; interpret the resulting interval in context.
    \end{itemize}
    \columnbreak
    \textbf{Tier E --- Extension}
    \begin{itemize}
        \item Complete all tiers plus Tier E.
        \item Compare two CIs and explain what each implies about the population means.
    \end{itemize}
\end{multicols}
\end{skillbox}

\begin{skillbox}[Individual Work \& Assessment]{redbox}
\begin{minipage}[t]{0.48\linewidth}
    \textbf{Exit Ticket (5 min)}\\
    A nutritionist wants to estimate the difference in mean daily sugar intake (g) between
    adults in two cities. City X: $n = 35$, $\bar x = 82$ g, $s = 14$ g.
    City Y: $n = 40$, $\bar x = 75$ g, $s = 11$ g. Conditions are met.
    \begin{enumerate}
        \item Identify the procedure; state the point estimate.
        \item Compute the standard error.
        \item Interpret a 95\% CI of $(2.0, 12.0)$ g in context.
    \end{enumerate}
\end{minipage}
\hfill
\begin{minipage}[t]{0.48\linewidth}
    \textbf{AP-Style Check}\\
    \textit{A two-sample $t$-interval for $\mu_1 - \mu_2$ is computed as $(1.4, 8.6)$.
    Which conclusion is best supported?}
    \begin{enumerate}[label=(\Alph*)]
        \item There is no convincing evidence that $\mu_1 \ne \mu_2$.
        \item The true difference is exactly 5.0.
        \item We are confident $\mu_1 > \mu_2$ because the interval is entirely positive.
              \textcolor{keyred}{\textbf{$\leftarrow$ correct}}
        \item We need a larger sample to make any conclusion.
    \end{enumerate}
\end{minipage}
\end{skillbox}

\begin{skillbox}[Reinforcement \& Extension]{goldbox}
\begin{minipage}[t]{0.48\linewidth}
    \textbf{Homework Overview}
    \begin{itemize}
        \item Four to five problems: identify procedure, verify conditions, compute SE
              and ME, calculate the interval, and interpret.
        \item One problem requires students to analyze whether 0 is in the interval
              and state the implication for the population means.
    \end{itemize}
\end{minipage}
\hfill
\begin{minipage}[t]{0.48\linewidth}
    \textbf{Extension / Preview}
    \begin{itemize}
        \item \textit{Extension:} Change the confidence level from 95\% to 90\% for
              a given problem. How does the interval width change? Why?
        \item \textit{Preview:} Lesson 7.7 extends this logic to \emph{significance
              tests} for $\mu_1 - \mu_2$ --- the two-sample $t$-test.
    \end{itemize}
\end{minipage}
\end{skillbox}

\end{document}
